\section{ProtNLM：从氨基酸序列生成自然语言功能描述}

上一章仍以序列表示和 masked prediction 为主，ProtNLM 把输出进一步改成自然语言。这个任务的价值在于数据库中的 protein function 本来就以自由文本存在；如果模型能从 sequence 生成可检索、可读的描述，就能帮助研究者在大量未充分注释的蛋白质中提出候选功能。

\subsection{为什么蛋白质注释是语言任务}

标题页给出 ProtNLM，随后两页说明需求：测序产生的蛋白质序列增长很快，人工实验注释却昂贵缓慢；UniProt 等数据库包含大量自由文本描述，但深度和可靠性不均。模型要学习的不是“给序列起一个好听标题”，而是把 sequence evidence 映射到功能、家族、位置和生物过程的语言表述。

\lecturefigure{slide-50-protnlm-title.jpg}{ProtNLM：Model-based Natural Language Protein Annotation}{官方录像恢复幻灯片 V050；Stanford CS25 Lecture 17}

\lecturefigure{slide-51-protein-function-demand.jpg}{海量蛋白质序列与功能注释需求}{官方录像恢复幻灯片 V051；Stanford CS25 Lecture 17}

\lecturefigure{slide-52-uniprot-free-text.jpg}{UniProt 中的蛋白质自由文本注释}{官方录像恢复幻灯片 V052；Stanford CS25 Lecture 17}

\begin{knowledgebox}{术语表：protein annotation 的证据层级}
\term{sequence annotation} 从氨基酸序列预测标签或描述；\term{curated annotation} 由专家结合实验和文献整理；\term{automatic annotation} 根据相似性或模型批量传播；\term{experimental validation} 通过生化、结构或细胞实验确认功能。ProtNLM 主要改善前两层之间的信息接口，不能把生成文本升级为实验事实。
\end{knowledgebox}

\subsection{Sequence-to-language：把两种词表连接起来}

Sequence-to-language 页把输入的 amino acid token 与输出的 natural-language token 放在同一 encoder--decoder 框架中；captioning 示例说明模型可以生成连贯功能描述。这里的关键不是语言流畅，而是描述是否包含与输入序列一致的 family、domain、localization 或 process 信息。

\lecturefigure{slide-53-sequence-to-language.jpg}{从蛋白质序列到自然语言的 encoder--decoder 任务}{官方录像恢复幻灯片 V053；Stanford CS25 Lecture 17}

\lecturefigure{slide-54-protein-captioning.jpg}{蛋白质 captioning 的输入与生成描述示例}{官方录像恢复幻灯片 V054；Stanford CS25 Lecture 17}

给定氨基酸序列 $S=(s_1,\ldots,s_n)$ 和目标描述 $Y=(y_1,\ldots,y_T)$，条件生成目标为
$$
\mathcal{L}_{\text{caption}}=-\sum_{t=1}^{T}\log p(y_t\mid S,y_{<t}).
$$
其中 $s_i$ 是第 $i$ 个 amino acid，$y_t$ 是第 $t$ 个自然语言 token，$n$ 与 $T$ 分别是输入和输出长度。这个目标奖励与参考文本一致的生成，但如果数据库标签本身缺失或噪声很大，模型也会继承不确定性。

\begin{warningbox}{生成得像，不等于注释得真}
自由文本目标会奖励常见措辞，模型可能依据家族先验生成听起来合理但未经实验确认的描述。评价必须区分 lexical similarity、专业实体正确性、功能层级一致性和新颖预测；任何自动描述都应保留 provenance 与不确定性。
\end{warningbox}

训练集切分在这里尤其敏感。如果高度同源的蛋白质同时出现在训练和测试中，模型可能通过近邻记忆复述家族描述，看起来像从 sequence 推断功能。更严格的评价会按 sequence identity 或 protein family 分组，把远缘家族留到测试集，并分别报告常见功能与稀有功能。生成文本还可被拆成结构化槽位，例如 family、catalytic activity、subcellular location 与 evidence qualifier；逐槽位核验比一个整体 BLEU 或 ROUGE 分数更接近科学用途。自然语言接口很方便，但越方便越需要把证据等级显式写进输出。

\subsection{T5 多任务训练与结果}

T5 页显示 ProtNLM 不只做单向 captioning，而是把 span corruption、sequence-to-text 与相关任务统一为 text-to-text 形式。多任务训练让模型同时学习蛋白质局部模式和语言表述。结果页总结模型在内部评价与人类判断上的改进；应理解为注释质量和覆盖的证据，而不是对所有新蛋白质功能的实验确认。

\lecturefigure{slide-55-protnlm-t5-tasks.jpg}{ProtNLM 的 T5 风格预训练与多任务目标}{官方录像恢复幻灯片 V055；Stanford CS25 Lecture 17}

\lecturefigure{slide-56-protnlm-results.jpg}{ProtNLM 的蛋白质描述生成结果}{官方录像恢复幻灯片 V056；Stanford CS25 Lecture 17}

\begin{knowledgebox}{读图：结果需要拆成“覆盖”和“可信度”}
自动模型可以快速为更多序列给出描述，这是 coverage 收益；描述与参考注释或专家判断的一致性属于 quality 收益。两者必须同时看：只扩大覆盖会批量传播错误，只追求极高 precision 又可能拒绝大量未知蛋白质。
\end{knowledgebox}

\subsection{CRISPR-Cas9 案例：生成描述如何服务发现}

前面已经说明模型怎样生成蛋白质描述，本节把接口放入真实发现流程。应用页展示 CRISPR-Cas9 相关蛋白质空间；研究者面对的并不是一个已知答案，而是大量可能具有相似机制、不同大小或不同 targeting 特性的候选。模型生成或组织自然语言注释，可以帮助搜索序列、比较家族和形成实验假设。正确工作流是“模型筛选与解释候选—研究者检查数据库与结构证据—设计实验—验证功能”，而不是“模型写出描述—功能已被发现”。

\lecturefigure{slide-57-crispr-cas9-application.jpg}{应用示例：寻找 CRISPR-Cas9 相关蛋白质}{官方录像恢复幻灯片 V057；Stanford CS25 Lecture 17}

\teachervoice{讲者用 CRISPR 案例说明自然语言接口的科学价值：研究者不必只在 embedding 空间里寻找近邻，而可以用功能描述组织候选。但课堂措辞仍是 annotation 与 discovery support，没有把生成句子当作 wet-lab proof。}

\begin{importantbox}{ProtNLM 的正确定位}
ProtNLM 把蛋白质序列和人类可读知识连接起来，使检索、归纳和候选生成更高效。它最像“科学信息接口”：输入是 sequence，输出是可读 hypothesis；真正的生物功能仍由数据库 provenance、文献与实验闭环决定。
\end{importantbox}

\subsection{本章小结}

ProtNLM 将蛋白质表示学习扩展为 sequence-to-language，通过 T5 风格多任务训练生成自由文本功能描述。它能扩大注释覆盖、改善知识检索并支持候选发现，但文本生成质量与生物功能真实性必须分开评价。

